\chapter{Implementation Code Samples}
\label{appendix:code}

This appendix provides key implementation code for the data structures described in this thesis. The complete source code is available in the accompanying repository.

\section{ResizableBitVector Implementation}
\label{sec:rbv-code}

The following Java code implements the extensible bitvector with the $f = n/w$ resize strategy described in Chapter~\ref{chap:extensible}.

\begin{lstlisting}[language=Java, caption={ResizableBitVector.java - Core implementation}, label={lst:rbv}]
package com.chp3.efano;

import it.unimi.dsi.bits.BitVector;
import it.unimi.dsi.bits.LongArrayBitVector;
import it.unimi.dsi.sux4j.bits.SimpleSelect;

/**
 * Resizable bit vector with f = n/w resize strategy.
 * Achieves o(n) space overhead while maintaining O(1)
 * amortised append time.
 */
public class ResizableBitVector {
    public long length = 0;
    public long capacity = 64;
    private LongArrayBitVector bv;

    public ResizableBitVector() {
        bv = LongArrayBitVector.getInstance(capacity);
    }

    /**
     * Append a bit to the bitvector.
     * Resizes by n/64 when full (f = n/w strategy).
     */
    public void append(int bit) {
        if (length == capacity) {
            // Resize by n/w = n/64 for 64-bit words
            long newCapacity = length + (length >> 6);
            if (newCapacity == capacity) {
                newCapacity = capacity + 1;
            }
            bv.length(newCapacity);
            capacity = newCapacity;
        }
        bv.set(length, bit == 1);
        length++;
    }

    public BitVector bitVector() {
        return bv.subVector(0, length);
    }

    public SimpleSelect buildSelectIndex() {
        return new SimpleSelect(bv.subVector(0, length));
    }
}
\end{lstlisting}

\section{SIMD Predecessor Search Implementation}
\label{sec:simd-code}

The following C++ code implements the improved Fredman method using AVX2 SIMD instructions, as described in Chapter~\ref{chap:predecessor}.

\begin{lstlisting}[language=C++, caption={Improved Fredman rank using XOR-MinPos}, label={lst:fredman}]
#include <immintrin.h>
#include <cstdint>

/**
 * Compute rank using SIMD XOR-MinPos method.
 * Exploits: min(x XOR S[i]) = element with longest
 * common prefix with x.
 */
template<int N>
int rankFWImproved(const uint16_t* S, uint16_t x,
                   int CT[][64]) {
    // Broadcast query to all SIMD lanes
    __m256i a = _mm256_set1_epi16(x);

    // Load set elements
    __m256i b = _mm256_loadu_si256(
        (const __m256i*)S);

    // Compute XOR differences
    __m256i c = _mm256_xor_si256(a, b);

    // Find minimum (element with longest LCP)
    // Note: _mm256_minpos_epu16 operates on 128-bit
    __m128i lo = _mm256_extracti128_si256(c, 0);
    __m128i hi = _mm256_extracti128_si256(c, 1);
    __m128i min_lo = _mm_minpos_epu16(lo);
    __m128i min_hi = _mm_minpos_epu16(hi);

    // Extract index and value
    int idx_lo = _mm_extract_epi16(min_lo, 1);
    int idx_hi = _mm_extract_epi16(min_hi, 1);
    int val_lo = _mm_extract_epi16(min_lo, 0);
    int val_hi = _mm_extract_epi16(min_hi, 0);

    int i = (val_lo <= val_hi) ? idx_lo : idx_hi + 8;

    // Compute LCP for correction
    int j = __builtin_clzll(S[i] ^ x);

    return CT[i][j];
}
\end{lstlisting}

\section{Ajtai Method with Lookup Tables}
\label{sec:ajtai-code}

\begin{lstlisting}[language=C++, caption={Ajtai method using Answer and Correction tables}, label={lst:ajtai}]
#include <immintrin.h>
#include <cstdint>
#include <vector>

/**
 * Ajtai's constant-time predecessor search.
 * Uses pre-computed Answer Table (AT) and
 * Correction Table (CT).
 */
class AjtaiPredecessor {
private:
    std::vector<int> AT;      // Answer Table
    std::vector<std::vector<int>> CT;  // Correction Table
    uint64_t mask;            // Significant bit mask
    int b;                    // Number of sig. bits

public:
    void build(const uint16_t* S, int k) {
        // Compute significant bit positions
        mask = computeMask(S, k);
        b = __builtin_popcountll(mask);

        // Build Answer Table: O(2^b) entries
        AT.resize(1 << b);
        for (int i = 0; i < (1 << b); i++) {
            AT[i] = blindSearch(i, S, k);
        }

        // Build Correction Table: k * 64 entries
        CT.resize(k, std::vector<int>(64));
        for (int i = 0; i < k; i++) {
            for (int j = 0; j < 64; j++) {
                CT[i][j] = computeCorrection(i, j, S, k);
            }
        }
    }

    int rank(const uint16_t* S, uint16_t x) {
        // Compress query key
        uint64_t y = _pext_u64(x, mask);

        // Lookup in Answer Table
        int i = AT[y];

        // Compute LCP for correction
        int j = __builtin_clzll(S[i] ^ x);

        // Return corrected result
        return CT[i][j];
    }
};
\end{lstlisting}

\section{Dynamic Prefix Sum Operations}
\label{sec:prefixsum-code}

\begin{lstlisting}[language=C++, caption={Dynamic prefix sum with lazy synchronisation}, label={lst:prefixsum}]
#include <immintrin.h>
#include <cstdint>

/**
 * Dynamic searchable prefix sum using SIMD.
 * Based on Dietz's lazy synchronisation technique.
 */
class DynamicPrefixSum {
private:
    alignas(32) uint16_t A[16];  // Original array
    alignas(32) uint16_t B[16];  // Prefix sums
    alignas(32) uint16_t C[16];  // Update buffer
    int counter = 0;
    int m;  // Max gap in B

public:
    uint64_t sum(int i) {
        // B[i] + sum of updates in C[0..i]
        __m256i cv = _mm256_loadu_si256(
            (const __m256i*)C);
        // Horizontal sum up to position i
        return B[i] + horizontalSum(cv, i);
    }

    void increment(int i) {
        // Update buffer C
        C[i]++;
        counter++;

        // Refresh B after m updates
        if (counter >= m) {
            refresh();
            counter = 0;
        }
    }

    int select(uint64_t j) {
        // Use SIMD XOR-MinPos from Chapter 5
        // to find position in compressed B
        return simdSelect(j);
    }

private:
    void refresh() {
        // Synchronise B with accumulated updates
        for (int i = 0; i < 16; i++) {
            B[i] += horizontalSum(C, i);
        }
        // Reset update buffer
        std::memset(C, 0, sizeof(C));
        // Recompute max gap
        m = computeMaxGap();
    }
};
\end{lstlisting}
